\documentclass[a4paper, pagesize=pdftex, pdftex, twoside, headsepline, index=totoc,toc=listof, fontsize=10pt, cleardoublepage=empty, headinclude, DIV=13, BCOR=13mm]{scrartcl}

\usepackage[ngerman]{babel}
\usepackage{scrtime} % Bestandteil von KOMA-Skript, ermoeglicht Zugriff auf Uhrzeit des Kompilierens 
\usepackage{scrpage2} % ermöglicht Bearbeiten von Kopf- und Fusszeilen (wie fancyhdr, nur optimiert auf KOMA-Skript, leich andere Syntax)
\usepackage[utf8]{inputenc} % Gibt an in welcher Textcodierung der Code verstandne werden soll
\usepackage{etex} % sehr technisch, ermöglicht LaTeX mehr Speicher zu belegen
\usepackage[T1]{fontenc} % auch sehr technisch; ist wichtig, um die Schriftarten richtig zu behandeln
\usepackage{textcomp} %verhindert ein paar Fehler bei den Fonts
\usepackage{amsmath} % Packet der American Mathematical Society, das viele Mathematik-Umgebungen und -Befehle definiert
\usepackage{amssymb} %zusätzliche Symbole
\usepackage{latexsym} % nochmal zusätzliche Symbole
\usepackage{stmaryrd} % nochmal mehr zusätzliche Symbole, u.a. Blitz für Widerspruchsbeweise ;)
\usepackage{nicefrac} % schräge Brüche, benutzte ich für Quotienvektorräume
\usepackage{paralist} % redefiniert alle Listenbefehle, sodass diese einen optionalen Parameter haben, der die Nummerierung angibt
\usepackage{dsfont} % Schriftart für N,Z,Q,R die ich momentan benutze (mittels \mathds{R} z.B)
\usepackage[pdftex]{graphicx} % Packet, dass das Einbinden von Grafiken aus Dateien ermöglicht
\usepackage{makeidx}% ermöglicht das automatische Anlegen eines Index 
\usepackage{extarrows}

\flushbottom
\usepackage[normalem]{ulem}
\setlength{\ULdepth}{1.8pt}

%--Indexverarbeitung
\newcommand{\bet}[1]{\uline{\textbf{#1}}} %Betonung von Text
\newcommand{\Index}[1]{\uline{\textbf{#1}}\index{#1}} % Befehl, der gleichzeitg das Argument hervorhebt und in den Index mitaufnimmt
\makeindex % startet das automatische Sammeln der Index-Einträge
% Ein kleiner Text am Anfang des Index
\setindexpreamble{{\noindent \itshape Die \emph{Seitenzahlen} sind mit Hyperlinks zu den entsprechenden Seiten versehen, also anklickbar!} \par \bigskip}
\renewcommand{\indexpagestyle}{scrheadings} % Seitenstil für den Index festlegen

%--Farbdefinitionen
\usepackage[usenames, table, x11names]{xcolor} %usenames und x11names, aktivieren viele Farben; siehe Dokumentation von xcolor
% Es lassen sich natürlich auch eigene Farben definieren (hier nur Graustufen)
\definecolor{dark_gray}{gray}{0.45}
\definecolor{light_gray}{gray}{0.7}

%--Zum Zeichnen (ich habe es jetzt mal mit aufgenommen, aber es ist eigentlich nochmal ein ganz anderes Thema, sodass ich da jetzt nicht viel zu sagen werde)
\usepackage{tikz} % TikZ steht übrigens für "TikZ ist kein Zeichenprogramm", ein rekursives Akronym ...
\tikzset{>=latex}
\usetikzlibrary{shapes,arrows}
\usetikzlibrary{calc}
\usetikzlibrary{decorations.pathreplacing}
% Hiermit kann man ganz leicht kommutative Diagramme zeichnen (deswegen auch "cd")
\usepackage{tikz-cd}

%--Marginnote, ermöglicht es kleine Notizen an neben den eigentlichen Textkörper zu setzten
\usepackage{marginnote}
\renewcommand*{\marginfont}{\color{Honeydew4} \footnotesize }

%--Schriftarten
\usepackage{lmodern} % neuere Version der Standard-LaTeX-Schriftarten
\renewcommand{\familydefault}{\sfdefault} %Standardschriftart auf die serifenlose Schriftart setzen

%--Hyperref; aktiviert Hyperlinks in der erzeugten PDF-Datei und definiert deren Aussehen
\usepackage[colorlinks, pdfpagelabels, pdfstartview=FitH, bookmarksopen=true, bookmarksnumbered=true,linkcolor=black,urlcolor=SkyBlue2, plainpages=false, hypertexnames=false, citecolor=black, hypertexnames=true]{hyperref}

%--Römische Zahlen
\newcommand{\RM}[1]{\MakeUppercase{\romannumeral #1{}}}



%-- Definitionen von weiteren Mathe-Befehlen, die dann das "richtige" Aussehen haben. Hier sind der Phantasie keine Grenzen gesetzt
\DeclareMathOperator{\id}{id} %identische Abbildung
\DeclareMathOperator{\End}{End} %Endomorphismen
\DeclareMathOperator{\rg}{rg} %Rang
\DeclareMathOperator{\diam}{diam} %Durchmesser
\DeclareMathOperator{\dist}{dist} %Distanz
\DeclareMathOperator{\grad}{grad} %Gradient
\DeclareMathOperator{\rot}{rot} %Rotation
\DeclareMathOperator{\hess}{Hess} %Hesse-Matrix


%--Skalarprodukt (cooler Befehl, den ich im Internet gefunden habe; benutzt TeX-Befehle)
\makeatletter
\newcommand{\sprod}[2]{\ensuremath{%
  \setbox0=\hbox{\ensuremath{#2}}
  \dimen@\ht0
  \advance\dimen@ by \dp0
  \left\langle \left.#1 \,\rule[-\dp0]{0pt}{\dimen@}\right|#2\right\rangle}}
\makeatother

%--Norm (auch aus dem Internet, wird auch auf der Beispielseite verwandt)
\newcommand{\norm}[2][\relax]{
\ifx#1\relax \ensuremath{\left\Vert#2\right\Vert}
\else \ensuremath{\left\Vert#2\right\Vert_{#1}}
\fi}


%--selbstgeschriebenen Befehle
%--Betrag
\newcommand{\abs}[1]{\ensuremath{\left\vert#1\right\vert}}

%--Umklammern mit passender Größe der Klammern
\newcommand{\enbrace}[1]{\ensuremath{\left( #1\right)}}

%--Mengen
\newcommand{\penbrace}[1]{\ensuremath{\left\{#1\right\}}}

%--Differential
\newcommand{\diff}[2]{\ensuremath{\frac{\partial #1}{\partial #2} }}





%--Konfiguration von scrheadings (hierzu findet man auch massig zusätzliche Infos im scrguide)
\setheadsepline{1pt}[\color{light_gray}] 
\pagestyle{scrheadings} %definiert welcher Seitenstil im Dokument normalerweise verwendet werden soll
\clearscrheadfoot % erstmal alles, was eventuell in Kopf- und Fusszeile steht, löschen

% ohead steth für "outer Head", auf einer linken Seite als für den linken Kopf, auf einer rechten für den rechten (es wird ja ein doppelseitiges Layout verwandt)
% durch den hfill-Befehl wird die allesdings untergraben, was aber von mir beabsichtig ist
\ohead{\footnotesize \sffamily \color{light_gray} \includegraphics[height=0.5 cm, keepaspectratio]{Bilder/Logo_WWU_Muenster_light_gray.jpg}
	\includegraphics[height=0.4cm, keepaspectratio]{Bilder/fb10logo.jpg}
	 Tobias Wedemeier -- Einführung in die Algebra
	\hfill { \tiny \sffamily \color{light_gray} Stand: \today \; \thistime[:]}}

% durch den optionalen Parameter (eckige Klammern) wird auch der plain-Pagestyle bearbeitet, damit auch auf leeren Seiten die Seitenzahl steht
\ofoot[{ \color{dark_gray} \LARGE \sffamily \thepage}]{{ \color{dark_gray} \LARGE \sffamily \thepage}} 

% hier wird die aktuelle section in den "inner foot" gesetzt
\automark{section}
\ifoot{ \color{dark_gray} \small \leftmark}

%--Metadaten, die durch den \maketitle gesetzt werden
\author{Tobias Wedemeier}
\titlehead{\includegraphics[width=7cm, keepaspectratio]{Bilder/Logo_WWU_Muenster_light_gray.jpg}
			\includegraphics[width=6cm, keepaspectratio]{Bilder/fb10logo.jpg}}
\title{Einführung in die Algebra}
\subtitle{}
\publishers{gelesen von \\ Prof. Dr. Kramer}

\newcommand{\zz}{$\mathrm{Z\kern-.3em\raise-0.5ex\hbox{Z}}$} % zu zeigen ZZ aus dem inet
\setlength{\parindent}{0pt}%absatz nicht einrücken
\newcommand{\lh}[1]{\langle #1 \rangle} %lineare Hülle

\begin{document}
\maketitle
\thispagestyle{empty}
\newpage

\pagenumbering{Roman}

\tableofcontents
\cleardoubleoddemptypage %sorgt dafür, dass alles folgende erst auf der nächsten freien "rechten" Seite steht

\pagenumbering{arabic}
\setcounter{page}{1}

\section{Elementare Gruppentheorie}
\label{sec:elementare_gruppentheorie}

\bet{Erinnerung:} eine \Index{Verknüpfung} auf einer nicht leeren Menge $X$ ist eine Abbildung
\begin{equation*}
\begin{aligned}
X\times X \to X , (x,y) \mapsto m(x,y).
\end{aligned}
\end{equation*}
Häufig schreibt man $m(x,y)= x \* y$ oder $ m(x,y) = x + y$, je nach Kontext. Die Schreibweise $m(x,y)=x+y$ wird eigentlich nur für kommutative Verknüpfungen benutzt, d.h. wenn $\forall x,y\in X$ gilt $m(x,y)=m(y,x)$.

\subsection{Def. Gruppe}
\label{sub: gef_gruppe}
Eine \Index{Gruppe} $(G,\cdot )$ besteht aus einer Verknüpfung $\cdot $ auf einer nicht leeren Menge $G$, mit folgenden Eigenschaften:
\begin{enumerate}[(G1)]
	\item Die Verknüpfung ist \uline{assoziativ}, d.h. $(x\cdot y)\cdot z = x \cdot (y \cdot z)$ gilt $\forall x,y,z \in G$.\\ 
	(Folglich darf man Klammern weglassen.)
	\item Es gibt ein \uline{neutrales Element} $e \in G$, d.h. es gilt $e\cdot x= x\cdot e= x \forall x\in G$
	\item Zu jedem $x\in G$ gibt es ein \uline{Inverses} $y \in G$, d.h. $xy=e=yx.$\\
	man schreibt dann auch $y=x^{-1}$ für das Inverse zu x.
\end{enumerate}
Fordert man von der Verknüpfung nur (G1) und (G2), so spricht man von einer Halbgruppe mit Eins oder einem \Index{Monoid}. Fordert man nur (G1), so spricht man von einer Halbgruppe.\\
% sub end

\subsection{Beispiel 1}
\label{sub:beispiel_1}
\begin{itemize}
	\item $(\mathds{Z}, +), (\mathds{Q}, +)$ sind kommutative Gruppen.
	\item $(\mathds{Z},\cdot), (\mathds{N},\cdot), (\mathds{N}, +)$ sind Monoide.
\end{itemize}
%sub end

\subsection{Beobachtungen}
\label{sub: beobachtungen}
\begin{enumerate}[a)]
	\item Das Neutraleelement (einer Verknüpfung) ist eindeutig bestimmt: sind $e,e'$ beides Neutralelemente, so folgt: $e=ee'=e'$
	\item Das Inverse zu $x$ ist eindeutig bestimmt:\\
	$xy=e=xy'=y'x \Rightarrow y'=y'e=y'xy=ey=y$
\end{enumerate}
%sub end

\subsection{Lemma 1 (Sparsame Definition von Gruppen)}
\label{sub:lemma_1}
Sei $G \times G \to G$ eine assoziative Verknüpfung. Dann ist $G$ schon eine Gruppe, wenn gilt:
\begin{enumerate}[(i)]
	\item es gibt $e \in G$ so, dass $ex=x$ $\forall x \in G$ gilt.
	\item zu jedem $x\in G$ gibt es ein $y \in G$ mit $ yx=e$
\end{enumerate}
\bet{Beweis}\\
Sei $yx=e$, es folgt $yxy=y$. Wähle $z$ mit $zy=e$, es folgt $\underbrace{zy}_{=e}xy=zy=e \Rightarrow xy=e$\\
Weiter gilt $xe=xyx=ex=x$.
\hfill $\square$
%sub end

\subsection{Beispiel 2}
\label{sub: beispiel_2}
Sei $X$ eine nicht leere Menge, sei $X^X=\{f : X \to X\}$ die Menge aller Abbildungen von $X$ nach $X$. Als Verknüpfung auf $X$ nehmen wir die Komposition von Abbildungen. Dann gilt wegen $f=id_X \circ f= f \circ id_X$, dass $id_X$ ein Neutralelement ist.\\
Damit haben wir ein Monoid $(X_X, \circ )$.\\
Sei $Sym(X)=\{f:X\to X | f$ bijektiv$\}$. Zu jedem $f\in Sym(X)$ gibt es also eine Umkehrabbildung $g:X\to X$ mit $f \circ g=g\circ f=id_X$. Folglich ist $(Sym(X), \circ)$ eine Gruppe, die \bet{Symmetrische Gruppe}\index{Gruppe!symmetrische !}. Wenn X endlich ist mit $n$ Elementen, so gibt es genau $n!=n(n-1)(n-2)\dots \cdot 2\cdot 1$ Permutationen, also hat $Sym(X)$ dann genau $n!$ Elemente.\\
Für $X=\{1,2,3,\dots,n\}$ schreibt man auch $Sym(X)=Sym(n)\bigg( =S_n \bigg)$.
%sub end

\subsection{Def. zentralisieren}
\label{sub:def_zentralisieren}
Sei $G \times G \to G$ eine Verknüpfung. Wir sagen, $x,y \in G$ vertauschen oder kommutieren oder $x$ \Index{zentralisiert} $y$, wenn gilt $xy=yx$.\\
Eine Gruppe, in der alle Elemente vertauschen heißt kommutativ oder \Index{abelsch}. 
%sub end

\subsection{Beispiel 3}
\label{sub: beispiel_3}
\begin{enumerate}[(a)]
	\item $(\mathds{Z}, +), (\mathds{Q}, +), (\mathds{Q}^*,\cdots)$ sind abelsche Gruppen.
	\item $K$ Körper, $G=Gl_2(K)=\{X \in K^{2\times 2}~ |~ det(X)\not= 0 \}$ Gruppe der invertierbaren $2 \times 2$ Matrizen.\\
\[
	\begin{pmatrix}	1 & 1\\ 0 & 1 \end{pmatrix}
	\begin{pmatrix} 0 & 1\\	1 & 0 \end{pmatrix}
	=
	\begin{pmatrix}	1 & 1\\	1 & 0 \end{pmatrix}
\]
\[
	\begin{pmatrix} 0 & 1\\	1 & 0 \end{pmatrix}
	\begin{pmatrix} 1 & 1\\ 0 & 1 \end{pmatrix}
	=
	\begin{pmatrix} 0 & 1\\ 1 & 1 \end{pmatrix}
\]
	$\Rightarrow$ nicht abelsch, genauso $Gl_n(K)$ für $n\ge 2$.
	\item $Sym(2)$ ist abelsch, aber $Sym(3)$ nicht. Allgemein ist $Sym(X)$ nicht abelsch, falls $\#X \ge 3$ gilt.
\end{enumerate}
%sub end

\subsection{Def. Untergruppe}
\label{sub: def_untergruppe}
Sei $G$ eine Gruppe, sei $H\subseteq G$. Wir nennen $H$ \bet{Untergruppe} \index{Gruppe! Unter-!} von $G$, wenn gilt:
\begin{enumerate}[(UG1)]
	\item $e\in H$
	\item $x,y\in H \Rightarrow xy\in H$
	\item $x\in H \Rightarrow x^{-1} \in H$
\end{enumerate}
Offensichtlich ist eine Untergruppe dann wieder eine Gruppe, mit der von $G$ vererbten Verknüpfung.

\minisec{Bsp}
\begin{enumerate}[(a)]
	\item $(\mathds{Q}, +)$. $\mathds{Z}$  ist Untergruppe, denn $0 \in \mathds{Z}, m,n \in \mathds{Z} \Rightarrow m+n\in \mathds{Z}$ und $n\in \mathds{Z} \Rightarrow -n\in \mathds{Z}$
	\item $(\mathds{Q}^*,\cdot)$. $\mathds{Z}^*$ ist keine Untergruppe, kein Inverses.
\end{enumerate}
%sub end

\subsection{Lemma 2}
\label{sub:lemma_2}
Sei $G$ eine Gruppe und sei $U$ eine nicht leere Menge von Untergruppen von $G$. Dann ist auch\\
$\bigcap U = \{g\in G | \forall H\in U$ gilt $g\in H\}$\\
eine Untergruppe von $G$.\\
\bet{Beweis}\\
Für alle $H\in U$ gilt $e\in H$, also $e\in \bigcap U$. Angenommen $x,y\in \bigcap U$. Dann gilt für alle $H\in U$, dass $xy\in H$ sowie $x^{-1}\in H$. Es folgt $xy\in \bigcap U$ sowie $x^{-1}\in \bigcap U$.
\hfill $\square$
%sub end

\subsection{Def. $\lh{X}$}
\label{sub:def_lhX}
Sei $G$ eine Gruppe und $X \subseteq G$ eine Teilmenge. Wir setzen:
\[\lh{X}=\bigcap\{H\subseteq G | H \text{ Untergruppe und } X \subseteq H\}\]
Ist nicht leer, da mindestens $G$ enthalten ist.
\begin{itemize}
	\item Es gilt z.B. $\lh{\emptyset}=\{e\}$, denn $\{e\}$ ist Untergruppe.
	\item Ist $H \subseteq G$ Untergruppe mit $X \subseteq H$, so folgt $X\subseteq \lh{X} \subseteq H$, insb. also $\lh{H}=H$.
\end{itemize}

\minisec{Satz}
Sei $X \subseteq G$ und sei $W=\{x_1\cdot x_2,\cdot \dots \cdot x_s | s\ge 1, x_i\in X \text{ oder } x_i^{-1}\in X ~\forall i=1,\dots,s\}$.\\
Dann gilt: $ \lh{X}=\{e\}\cup W$.\\
\bet{Beweis}\\
Wegen $X\subseteq \lh{X}$ und $e\in \lh{X}$ folgt $\{e\}\cup W\subseteq \lh{X}$. Ist $f,g\in W$, so folgt $fg\in W$ sowie $f^{-1}\in W$, also ist $H=\{e\}\cup W$ eine Untergruppe von $G$, mit $X\subseteq H$. Es folgt $\lh{X}\subseteq H=\{e\}\cup W$.
\hfill $\square$
%sub end

\subsection{Def. zyklische Gruppe}
\label{sub: def_zklische_gruppen}
Sei $G$ eine Gruppe und sei $g\in G$. Für $n\ge 1$ setze $g^n=\underbrace{g\cdot \dots \cdot g}_{n-mal}$ sowie $g^{-n}=\underbrace{g^{-1}\cdot \dots \cdot g^{-1}}_{n-mal}$ und $g^0=e$.\\
Dann gilt $\forall k,l\in \mathds{Z}$, dass $g^k\cdot g^l = g^{k+l}$.\\
Sei $\lh{g}=\lh{\{g\}}\stackrel{1.10}{=}\{g^n | n\in \mathds{Z}\}$. Man nennt $\lh{g}$ die von $g$ erzeugte \bet{zyklische Gruppe}\index{Gruppe!zyklische !}. Wenn für ein $n\ge 1$ gilt $g^n=e$, so heißt $n$ ein \Index{Eponent} von $g$. DIe \Index{Ordnung} von g ist der kleinste Eponent von g,
\[o(g)=min\enbrace{\{n\ge 1 | g^n=1\}\cup \{\infty\}}\]
$o(g)=\infty$ bedeutet: $g^n\not= e~\forall n\ge 1$\\
$o(g)=1$ bedeutet: $g^n=g=e$
%sub end

\subsection{Zyklische Gruppen}
\label{sub:zyklische_gruppen}
Eine Gruppe $G$ heißt \Index{zyklisch}, wenn es ein $g\in G$ gibt mit $G=\lh{g}$. Wegen $g^kg^l=g^{k+l}=g^{l+k}=g^lg^k$ gilt: zyklische Gruppen sind abelsch.

\minisec{Satz}
Sei $G=\lh{g}$ zyklisch mit $o(g)=n<\infty$. Dann gilt $\#G=n$ und $G=\{g,g^1,g^2,g^3,\dots,g^n\}$.
\bet{Beweis}\\
Jedes $m\in \mathds{Z}$ lässt sich schreiben als $m=kn+l$ mit $0\le l<n$ (Teilen mit Rest), also $g^m= \underbrace{g^{kn}}_{=e}.g^l=g^l$. Es folgt $G\subseteq \{g,g^2,\dots,g^n\}, g^n=g^0$.
Ist $g^k=g^l$ für $0\le k\le l<n$, so gilt $e=g^0=g^{l-k}$, also $l-k=0$ (wegen$l<n$), also $\#\{g,g^2,\dots,g^n=g^0\}=n$.
\hfill $\square$

\minisec{Folgerung}
Ist $G$ endlich mit $\#G=n$ und ist $h\in G$ mit $O(h)=n$, so folgt $\lh{h}=G$. Insbesondere ist dann $G$ eine zyklische Gruppe.
\hfill $\square$
%sub end

\subsection{Nebenklassen}
\label{sub:nebenklassen}
\index{Nebenklassen!Links-} \index{Nebenklassen!Rechts-}

Sei $G$ eine Gruppe und sei $H$ eine Untergruppe. Sei $a\in G$. Wir definieren:
\[aH=\{ah | h\in H\}\subseteq G\]
\[Ha=\{ha | h\in H\}\subseteq G\]
Man nennt $aH$ die \bet{Linksnebenklassen} von $a$ bzgl. $H$ (und $Ha$ die \bet{Rechtsnebenklassen}). In nicht abelschen Gruppen gilt im allgemeinen $aH\not=Ha$.

\minisec{Lemma}
Sei $H\subseteq G$ Untergruppe der Gruppe $G$ und $a,b\in G$.\\
Dann sind äquivalent:
\begin{enumerate}[(i)]
	\item $b\in aH$
	\item $bH=aH$
	\item $bH \cap aH \not= \emptyset$
\end{enumerate}
\bet{Beweis}\\
\begin{itemize}
	\item[$(i)\Rightarrow (ii):~b\in aH \Rightarrow b=ah$ für ein $h\in H \Rightarrow bH=\{ahh' | h' \in H\} \stackrel{H\text{ Untergruppe}}{=}\{ah'' | h''\in H\}=aH$]
	\item[$(ii) \Rightarrow (iii):$ klar]
	\item[$(iii) \Rightarrow (i):$ Sei $g \in bH \cap aH,~g=bh=ah' \Rightarrow b=ah'h^{-1} \in aH$, da $H$ Untergruppe]
\end{itemize}
\hfill $\square$

\minisec{Folgerung}
Jedes $g\in G$ liegt in genau einer Linksnebenklasse bzgl. $H$, nämlich $g \in gH$.
Entsprechendes gilt natürlich für Rechtsnebenklassen. Man setzt:\\
$\nicefrac{G}{H}=\{gH~|~g \in G \}$ Menge der Linksnebenklasse, Rechtsnebenklassen analog.

\minisec{Lemma}
Sei $H\subseteq G$ Untergruppe der Gruppe $G$, sei $a \in G$.\\
Dann ist die Abbildung $H \to gH, h \mapsto gH$ bijektiv.\\
\bet{Beweis}\\
"Surjektiv" ist klar nach Definition von $gH$. Angenommen, $gh=gh' \Rightarrow h=g^{-1}gh'=h'$
\hfill $\square$
%sub end

\subsection{Satz von Lagrange}
\label{sub:satz_von_lagrange}
\index{Satz von Lagrange}
Sei $G$ eine Gruppe und $H\subseteq G$ eine Untergruppe. Wenn zwei der drei Mengen $G,H,\nicefrac{G}{H}$ endlich sind, dann ist die dritte ebenfalls endlich und es gilt:\\
\[\#G=\#h \cdot \#\nicefrac{G}{H} \]
Insbesondere ist dann $\#H$ eine \Index{Teiler} von $\#G$.\\
\vfill
\bet{Beweis}\\
Wenn $G$ endlich ist, dann sind auch $H$ und $\nicefrac{G}{H}$ endlich.\\
Angenommen, $\nicefrac{G}{H}$ und $H$ sind endlich. Dann ist auch $G= \bigcup \nicefrac{G}{H}=\bigcup\{gH~|~gH\in \nicefrac{G}{H} \}$ endlich, da $\#gH=\#H$ nach \hyperref[sub:nebenklassen]{1.13}.\\
Jetzt zählen wir genauer: sei $\#\nicefrac{G}{H}=m; \#H=n$ etwa $\nicefrac{G}{H}=\{g_1H,g_2H,\dots g_mH \}$.\\
$g_iH\stackrel{1.13}{=}n~~~~~g_iH\cap g_jH=\emptyset$ für $i\not=j$ nach 1.13.\\
$G=g_1\cap \#g_2H\cap \dots \cap g_mH \Rightarrow \#G=m\cdot n$
\hfill $\square$

\minisec{Bem}
\begin{enumerate}[(1)]
	\item Eine entsprechende Aussage gilt für Rechtsnebenklassen.
	\item Die Abbildung $G \to G,~~g\mapsto g^{-1}$ bildet die Linksnebenklassen bijektiv auf die Rechtsnebenklassen ab:
	\[
	(gH)^{-1}=\{(gh)^{-1}~|~h \in H \} \stackrel{\text{Achtung!}}{=}\{h^{-1}g^{-1}~|~h \in H \}=\{hg^{-1}~|~h\in H \}=Hg^{-1}~~~~~~~~\text{(ÜA)}
	\]
\end{enumerate}

\minisec{Korollar A (Lagrange)}
Sei $G$ eine endliche Gruppe und sei $g\in G$. Dann teilt $o(g)$ die Zahl $\#G$.\\
\bet{Beweis}\\
Da $G$ endlich ist, folgt $o(g)<\infty$. Nach dem \hyperref[sub:satz_von_lagrange]{Satz von Lagrange} ist $\#\lh{g}=o(g)$ ein Teiler von $\#G$.
\hfill $\square$

\minisec{Korollar B}
Sei $G$ eine endliche Gruppe, sei $p$ eine \Index{Primzahl}  (d.h. die einzigen Teiler von $p$ sind 1 und $p$) und $p>1$. Wenn gilt $\#G=p$, dann ist $G$ zyklisch. Für jedes $g\in G \textbackslash\{e\}$ gilt $\lh{g}=G$.\\
\bet{Beweis}\\
Sei $g \in G\textbackslash\{e\}$. Dann ist $o(g)>1$ und $o(g)$ teilt $p$. Es folgt $o(g)=p$, also $G=\lh{g}$ vgl. \hyperref[sub:zyklische_gruppen]{1.12}.
\hfill $\square$

Für endliche Gruppen sind Teilbarkeitseigenschaften wichtig, wie wir sehen werden.\\
Die Zahl $\#\nicefrac{G}{H}:=[G:H]$ nennt man auch den \Index{Index von H in G}.

\minisec{Wichtige Rechenregeln in Gruppen}
\begin{enumerate}[(a)]
	\item Man darf \uline{kürzen}
	\begin{equation*}
	\begin{aligned}
		ax &= ay \Rightarrow x=y\\
		xa &= ya \Rightarrow x=y
	\end{aligned}
	\end{equation*}
	(multipliziere beide Seiten von links/rechts mit $a^{-1}$)
	\item Es gilt $(x^{-1})^{-1}=x$   ($x^{-1}x=e=xx^{-1} \Rightarrow (x^{-1})^{-1}=x$)
	\item Beim Invertieren darf die Reihenfolge umgedreht werden:\\
	\[(ab)^{-1}=b^{-1}a^{-1}\\
	\enbrace{ab(b^{-1}a^{-1})=e=(b^{-1}a^{-1})ab \Rightarrow (ab)^{-1}=b^{-1}a^{-1}}\\
	(\text{in abelschen Gruppen gilt natürlich damit} (ab)^{-1}=a^{-1}b^{-1} )
	\]
\end{enumerate}
% sub end
\subsection{Homomorphismen}
\label{sub:homomorphismen}
Seien $G,K$ Gruppen. Eine Abbildung $\varphi: G \to K$ heißt \bet{(Gruppen-)Homomorphismus}\index{Homomorphismus!Gruppen-}, wenn $\forall x,y \in G$ gilt
\[\varphi\underbrace{(x\cdot y)}_{\text{Verküpfung in G} } =\underbrace{\varphi(x)\varphi(y)}_{\text{Verknüpfung in K}} \]

\minisec{Bsp}
\begin{enumerate}[(a)]
	\item $id_G: G \to G$ ist Homomorphismus
	\item $H \subseteq G$ Untergruppe   $i:H \hookrightarrow G,~~~h \mapsto h$ Inklusion, ist Homomorphismus.
	\item $(G,\cdot)=(\mathds{Z},+)~~m\in \mathds{Z} ~~ \varphi:\mathds{Z} \to \mathds{Z}, x\mapsto mx$ ist Homomorphismus, denn $\phi(x+y)=m(x+y)=mx+my=\varphi(x)+\varphi(y)$
	\item $G$ Gruppe, $a \in G,~ a\not= e,~~ \lambda_a(x)=ax$.\\
	$\lambda: G \to G$ ist kein Homomorphismus, denn $\lambda_a(e)=a, \lambda(ee)=a$, aber $\lambda_a(e)\lambda_a(e)=aa\not=a$
\end{enumerate}

\minisec{Lemma}
Sei $\varphi:G \to K$ ein Homomorphismus von Gruppen. Dann gilt $\varphi(e_G)=e_K$ und $\varphi(x^{-1})=\varphi(x)^{-1}~\forall x \in G$. ($e_G$ Neutralelement in $G$ und $e_K$ Neutralelement in $K$)\\
\bet{Beweis}\\
\[
	\varphi(e_G)=\varphi(e_G \cdot e_G)=\varphi(e_G) \cdot \varphi(e_G)
	\stackrel{\text{kürzen}}{\Rightarrow} e_K=\varphi(e_G)
\]
\[
	e_K=\varphi(e_G)=\varphi(x^{-1}x)=\varphi(x^{-1})\varphi(x) \Rightarrow \varphi(x)^{-1}=\varphi(x^{-1})
\]
\hfill $\square$

\uline{Achtung:} $\varphi(x)^{-1}$ ist das Inverse in $K$ von $\varphi(x)$ \uline{nicht} die Umkehrabbildung!\\

Das \Index{Bild} eines Homomorphismus $\varphi:G \to K$ ist $\varphi(G)\subseteq K$,\\
der \Index{Kern} ist $ker(\varphi)=\{x \in G~|~\varphi(x)=e_K \}\subseteq G$
%sub end

\subsection{Satz 1}
\label{sub:satz_1}
Bild und Kern von Gruppenhomomorphismen sind Untergruppen.\\

\bet{Beweis}\\
Setze $H=\varphi(G)\subseteq K$. Es folgt $e_K \in H$. Für $\varphi(x),\varphi(y)\in H$ gilt $\varphi(x)\varphi(y)=\varphi(xy)\in H$ sowie $\varphi(x)^{-1}=\varphi(x^{-1}) \in H$, also ist $H$ Untergruppe.




%sec end













\cleardoubleoddemptypage
\pagenumbering{Alph}
\setcounter{page}{1}
\printindex
\listoffigures
\end{document}